\section{Математическая постановка задачи}

В данном разделе задача верификации соответствия формальной спецификации и
исходного кода формализуется в терминах системы переходов, трассы наблюдений и
предиката согласованности шага. На этой основе задаётся строгий критерий
согласованности трассы с моделью как условие существования допустимого
поведения модели, совместимого с наблюдаемой трассой.

\subsection{Формализация модели и ее поведения}

Для формализации модели (спецификации) программы используется понятие
\emph{системы переходов} (transition system) \cite{Keller1976} с выделенным
множеством начальных состояний.

\begin{definition}
Системой переходов называется тройка
\begin{equation}\label{f1}
\mathcal{M} = (X, X_0, R),
\end{equation}
где
\begin{itemize}
    \item $X$ --- множество состояний;
    \item $X_0 \subseteq X$ --- множество начальных состояний;
    \item $R \subseteq X \times X$ --- отношение переходов.
\end{itemize}
\end{definition}

Элемент $x \in X$ интерпретируется как полное абстрактное описание состояния
системы в некоторый момент её выполнения. Условие $(x,y) \in R$ означает, что
из состояния $x$ система может перейти в состояние $y$ за один допустимый шаг.
Иначе говоря, отношение $R$ задаёт все допустимые одношаговые изменения
состояния системы.

Следовательно, система переходов $\mathcal{M}$ задаёт не одно конкретное
выполнение, а всё множество возможных выполнений модели. Именно это множество
допустимых выполнений и используется далее как эталонное описание поведения
системы.

\begin{definition}
Конечным поведением системы переходов $\mathcal{M}$ длины $n \in \mathbb{N}_0$
называется последовательность состояний
\begin{equation}\label{f2}
\pi_n = (x_0, x_1, \dots, x_n) \in X^{n+1},
\end{equation}
для которой выполнены условия
\[
x_0 \in X_0,
\qquad
\forall i \in \{0,\dots,n-1\} : (x_i,x_{i+1}) \in R.
\]
\end{definition}

Первое условие означает, что поведение начинается из некоторого допустимого
начального состояния. Второе условие означает, что каждый следующий элемент
последовательности получается из предыдущего допустимым переходом системы.
Таким образом, конечное поведение есть путь в ориентированном графе,
порождаемом системой переходов, начинающийся в одной из вершин из множества
$X_0$.

\begin{definition}
Для каждого $n \in \mathbb{N}_0$ обозначим через $\mathrm{B}_n(\mathcal{M})$
множество всех конечных поведений системы переходов $\mathcal{M}$ длины $n$:
\[
\mathrm{B}_n(\mathcal{M})
=
\left\{
(x_0,\dots,x_n) \in X^{n+1}
:
x_0 \in X_0,\;
\forall i \in \{0,\dots,n-1\} : (x_i,x_{i+1}) \in R
\right\}.
\]
Тогда множество всех конечных поведений системы переходов $\mathcal{M}$
обозначается через $\mathrm{B}(\mathcal{M})$ и определяется равенством
\begin{equation}\label{f3}
\mathrm{B}(\mathcal{M})
=
\bigcup_{n \in \mathbb{N}_0} \mathrm{B}_n(\mathcal{M}).
\end{equation}
\end{definition}

Множество $\mathrm{B}(\mathcal{M})$ содержит все конечные пути в системе
переходов, начинающиеся из допустимых начальных состояний. Тем самым оно
полностью определяет допустимое наблюдаемое поведение модели на конечных
префиксах.

Если некоторое исполнение реализации не может быть согласовано ни с одним
элементом множества $\mathrm{B}(\mathcal{M})$, то это означает, что наблюдаемое
поведение не принадлежит множеству допустимых поведений, задаваемых абстрактной
моделью. Следовательно, в этом случае фиксируется расхождение между реализацией
и спецификацией.

\subsection{Формализация трассы программы}

Реальная программа в процессе выполнения порождает наблюдаемую трассу. В
рассматриваемой постановке предполагается, что трасса не содержит отдельного
наблюдения начального состояния системы. Начальное состояние выбирается из
множества начальных состояний модели, а трасса фиксирует только наблюдения,
относящиеся к последующим шагам выполнения программы.

При этом в рамках рассматриваемой постановки предполагается, что каждой записи
трассы соответствует ровно один переход поведения модели. Иначе говоря, трасса
и поведение модели сопоставляются \emph{пошагово}.

\begin{definition}
Пусть $\Sigma$ --- множество всех возможных наблюдений. Тогда наблюдаемая
трасса длины $m \in \mathbb{N}_0$ задаётся последовательностью
\begin{equation}\label{f4}
\tau = (\sigma_1,\sigma_2,\dots,\sigma_m) \in \Sigma^m.
\end{equation}
\end{definition}

Каждый элемент $\sigma_i$ интерпретируется как наблюдение, относящееся к $i$-му
шагу выполнения программы. Таким образом, поведение модели задаётся
последовательностью состояний $(x_0,x_1,\dots,x_m)$, тогда как наблюдаемая
трасса задаётся последовательностью наблюдений
$(\sigma_1,\sigma_2,\dots,\sigma_m)$, в которой каждый элемент $\sigma_i$
относится к переходу из состояния $x_{i-1}$ в $x_i$.

\subsection{Согласованность трассы с моделью}

После введения понятий модели и трассы можно строго определить, что означает их
согласованность.

В общем случае трасса программы представляет собой последовательность неполных
изменений состояния системы: обычно она содержит лишь частичную информацию о
каждом шаге выполнения, существенную для верификации, например сведения о
произошедшем событии и об изменениях отдельных компонент состояния. Поэтому
необходимо формально задать, когда переход считается согласованным с данным
наблюдением.

\begin{definition}
Предикатом согласованности наблюдаемого шага называется отображение
\begin{equation}\label{f5}
C : \Sigma \times X \times X \to \{0,1\}.
\end{equation}
Значение $C(\sigma,x,y)=1$ означает, что наблюдение $\sigma$ согласованно с
переходом из состояния $x$ в состояние $y$.
\end{definition}

Предикат $C$ задаёт способ интерпретации трассы в терминах абстрактной модели.
Его конкретный вид определяется составом данных, записываемых в трассу.
Например, если наблюдение фиксирует имя произошедшего события и значения
некоторых изменяемых компонент состояния, то согласованность означает, что
рассматриваемый переход принадлежит классу переходов, соответствующему данному
событию, и приводит к зафиксированным изменениям состояния. Если же наблюдение
фиксирует значения лишь некоторых компонент следующего состояния, то
согласованность означает совпадение именно этих компонент, тогда как остальные
компоненты могут выбираться произвольно при условии допустимости перехода.

Таким образом, трасса задаёт не одно конкретное поведение системы, а систему
ограничений на каждый последующий переход.

\begin{definition}
Пусть дана трасса $\tau = (\sigma_1,\sigma_2,\dots,\sigma_m) \in \Sigma^m$.
Тогда поведение модели $\pi = (x_0,x_1,\dots,x_m) \in X^{m+1}$ называется
\emph{согласованным с трассой $\tau$}, если выполнены условия
\[
x_0 \in X_0,
\qquad
\forall i \in \{1,\dots,m\} : C(\sigma_i,x_{i-1},x_i)=1.
\]
\end{definition}

Иначе говоря, трасса длины $m$ задаёт ограничения на поведение модели,
состоящее из $m$ переходов и $m+1$ состояний. В рамках данной постановки не
рассматриваются сопоставления, при которых одному элементу трассы соответствует
несколько переходов модели или, наоборот, одному переходу модели соответствует
несколько элементов трассы. Такое ограничение соответствует используемой схеме
проверки, в которой каждая запись трассы обрабатывается за один шаг
ограниченной спецификации.

\begin{definition}
Множество всех последовательностей состояний (поведений), согласованных с
трассой наблюдений $\tau = (\sigma_1,\dots,\sigma_m)$, обозначается через
$\mathrm{M}(\tau)$ и определяется равенством
\begin{equation}\label{f6}
\mathrm{M}(\tau)
=
\left\{
(x_0,\dots,x_m) \in X^{m+1}
:
x_0 \in X_0,\;
\forall i \in \{1,\dots,m\} : C(\sigma_i,x_{i-1},x_i)=1
\right\}.
\end{equation}
\end{definition}

Важно подчеркнуть, что элементы множества $\mathrm{M}(\tau)$ сами по себе
ещё не обязаны быть допустимыми поведениями модели. В этом множестве
учитываются только ограничения, задаваемые трассой. Чтобы последовательность
состояний представляла собой корректное абстрактное поведение системы, она
дополнительно должна принадлежать множеству $\mathrm{B}(\mathcal{M})$.

Следовательно, трасса является согласованной с моделью тогда и только тогда,
когда существует хотя бы одно поведение, которое одновременно:
\begin{itemize}
    \item является допустимым поведением системы переходов;
    \item согласовано с наблюдаемой трассой.
\end{itemize}

Это приводит к следующему строгому определению.

\begin{definition}
Трасса $\tau$ называется \emph{согласованной} с моделью $\mathcal{M}$, если
\begin{equation}\label{f7}
\mathrm{B}(\mathcal{M}) \cap \mathrm{M}(\tau) \neq \varnothing.
\end{equation}
\end{definition}

Если же $\mathrm{B}(\mathcal{M}) \cap \mathrm{M}(\tau) = \varnothing$,
то ни одно допустимое поведение модели не может объяснить наблюдаемую трассу,
и, следовательно, фиксируется расхождение между реализацией и абстрактной
моделью.

При этом в данной постановке нельзя требовать выполнения включения
$\mathrm{M}(\tau) \subseteq \mathrm{B}(\mathcal{M})$. Действительно, множество
$\mathrm{M}(\tau)$ определяется только наблюдаемой трассой и, вообще говоря,
строится по неполной информации о выполнении программы. Поэтому оно содержит
все последовательности состояний, согласованные с зафиксированными
наблюдениями, включая и такие последовательности, которые удовлетворяют этим
наблюдениям, но не являются допустимыми поведениями модели. Иными словами, из
согласованности последовательности с трассой ещё не следует, что эта
последовательность принадлежит множеству $\mathrm{B}(\mathcal{M})$.

\newpage
\subsection{Итоговая математическая постановка задачи}

Пусть заданы:
\begin{itemize}
    \item модель - система переходов $\mathcal{M} = (X, X_0, R)$, где $X$ ---
множество состояний, $X_0 \subseteq X$ --- множество начальных состояний, $R
\subseteq X \times X$ --- отношение переходов;

    \item трасса программы $\tau = (\sigma_1,\sigma_2,\dots,\sigma_m) \in
\Sigma^m$, где $\Sigma$ - множество наблюдений

    \item предикат согласованности шага
$C : \Sigma \times X \times X \to \{0,1\}$.
\end{itemize}

Тогда задача верификации соответствия трассы программы и модели состоит в том,
чтобы определить, существует ли последовательность состояний
$(x_0,x_1,\dots,x_m) \in X^{m+1}$, для которой выполнены условия $x_0 \in X_0,
\forall i \in \{1,\dots,m\} : C(\sigma_i,x_{i-1},x_i)=1$.

Иными словами, требуется установить, существует ли хотя бы одно допустимое
поведение модели, пошагово согласованное с наблюдаемой трассой программы:
\begin{equation}\label{f8}
\left\{
(x_0,\dots,x_m) \in X^{m+1}
:
x_0 \in X_0,\;
\forall i \in \{1,\dots,m\} : C(\sigma_i,x_{i-1},x_i)=1
\right\}.
\end{equation}

Эквивалентно, требуется установить непустоту пересечения множествa , то есть проверить истинность условия
\begin{equation}\label{f9}
\mathrm{B}(\mathcal{M}) \cap \mathrm{M}(\tau) \neq \varnothing,
\end{equation}
где $\mathrm{B}(\mathcal{M})$ - множество всех конечных поведений системы
переходов $\mathcal{M}$, а $\mathrm{M}(\tau)$ - множество всех конечных
поведений, согласованных с трассой $\tau$.

Следовательно, рассматриваемая задача является \emph{задачей разрешимости}
(decision problem): по заданным объектам $\mathcal{M}, \tau, C$ необходимо
выдать бинарный ответ на вопрос, истинно или ложно утверждение \eqref{f8}.
